\addcontentsline{toc}{section}{Abstract}
\section*{Abstract}


Proximity-induced superconducting electromigrated break-junctions (SEMBJs) are highly specialized tools with unrivaled capabilities in molecular electronics, such as their ability to probe the magnetic edge states of graphene, for example.

This thesis investigates the design and fabrication of bowtie geometry palladium-niobium and gold-niobium based proximity-induced SEMBJs using double angle shadow mask evaporation and feedback-controlled electromigration. The effect of the dimensions of the device on its effective energy gap, which determines the temperatures and bias voltages at which the device maintains its superconducting characteristics, was characterized and optimized. This was done using low-temperature transport measurements and SEM imaging.

This thesis shows that creating a proximity induced SEMBJ using this method is feasible. Optimized design parameters were found with which the device has a proximity induced energy gap of 0.28 meV at the tips of the leads of the SEMBJ. This gap size is comparable to values reported in previous studies, indicating that the device design and fabrication method are effective for achieving high-quality superconductivity in this type of junction.

This thesis demonstrates that controlling the location of the electromigration induced gap is crucial to ensure a large energy gap. We further demonstrate that bowtie constricted gold films with a thickness of 15 nm can be successfully electromigrated with a high degree of control, and that 30 nm thick gold films can be effectively proximitized by a 50 nm thick niobium layer. These findings can be used to guide the design of SEMBJs used in future research. 

The devices show an asymmetric critical current upon applying a magnetic field in the perpendicular in-plane direction. Our device offers a simple lithographic platform for making Josephson-diode based devices.
















\pagenumbering{roman}

\newpage